%======================================================================
% HABILIDAD: Panel de Control / Dashboard
% ID: dashboard
% Descripción: Interfaz y funcionalidades del panel principal
%======================================================================

\capitulo{Panel de Control}

\section{Visión General del Dashboard}

El Panel de Control es la interfaz principal desde donde podrá acceder a todas las funcionalidades del sistema. Está diseñado para proporcionar una vista rápida del estado del sistema y accesos directos a las operaciones más frecuentes.

\subsection{Elementos Principales}

\begin{description}
    \item[Barra Superior] Contiene el menú de navegación, notificaciones y perfil de usuario.
    \item[Barra Lateral] Proporciona acceso a los diferentes módulos y secciones.
    \item[Área de Contenido] Espacio principal donde se muestra la información y herramientas activas.
    \item[Estado del Sistema] Indicadores visuales del estado general del sistema.
\end{description}

\section{Navegación}

\subsection{Menú Principal}

El menú principal está organizado en las siguientes categorías:

\begin{itemize}
    \item \faIcon{home} \textbf{Inicio} - Vista general y accesos rápidos
    \item \faIcon{th-large} \textbf{Aplicaciones} - Gestión de aplicaciones y servicios
    \item \faIcon{server} \textbf{Recursos} - Monitorización de recursos del sistema
    \item \faIcon{shield-alt} \textbf{Seguridad} - Configuraciones de seguridad y acceso
    \item \faIcon{users} \textbf{Usuarios} - Administración de usuarios y permisos
    \item \faIcon{sliders-h} \textbf{Configuración} - Ajustes del sistema
    \item \faIcon{chart-line} \textbf{Reportes} - Análisis y estadísticas
    \item \faIcon{question-circle} \textbf{Ayuda} - Documentación y soporte
\end{itemize}

\subsection{Accesos Directos}

En la sección de accesos directos podrá personalizar los enlaces a las funciones más utilizadas:

\begin{enumerate}
    \item Haga clic en el icono \faIcon{plus} para agregar un nuevo acceso directo
    \item Seleccione la función deseada del listado
    \item Arrastre para reordenar según sus preferencias
    \item Haga clic en \faIcon{trash} para eliminar un acceso directo
\end{enumerate}

\section{Widgets del Dashboard}

El dashboard puede personalizarse con diferentes widgets informativos:

\subsection{Estadísticas en Tiempo Real}

Muestra métricas clave del sistema:
\begin{itemize}
    \item Uso de CPU
    \item Consumo de memoria
    \item Espacio en disco
    \item Conexiones activas
\end{itemize}

\subsection{Actividad Reciente}

Lista los eventos más recientes del sistema:
\begin{itemize}
    \item Inicios de sesión
    \item Cambios de configuración
    \item Alertas y notificaciones
    \item Tareas completadas
\end{itemize}

\subsection{Estado de Servicios}

Indicadores visuales del estado de los servicios principales:
\begin{itemize}
    \item \faIcon{check-circle} \textcolor{sandraGreen}{Servicio funcionando correctamente}
    \item \faIcon{exclamation-triangle} \textcolor{sandraAccent}{Advertencia}
    \item \faIcon{times-circle} \textcolor{red}{Servicio detenido o con error}
\end{itemize}

\section{Personalización del Dashboard}

\subsection{Organización de Widgets}

Puede personalizar la disposición de los widgets:

\begin{enumerate}
    \item Haga clic en \textbf{Personalizar Dashboard}
    \item Arrastre los widgets para reorganizarlos
    \item Redimensione los widgets arrastrando sus bordes
    \item Haga clic en \faIcon{times} para eliminar widgets
    \item Guarde la configuración
\end{enumerate}

\subsection{Temas y Apariencia}

El sistema soporta diferentes temas visuales:

\begin{itemize}
    \item \textbf{Tema Claro} - Interfaz con fondo blanco
    \item \textbf{Tema Oscuro} - Interfaz con fondo oscuro
    \item \textbf{Tema Automático} - Cambia según la configuración del sistema
\end{itemize}

Para cambiar el tema:
\begin{enumerate}
    \item Vaya a \textbf{Configuración > Apariencia}
    \item Seleccione el tema deseado
    \item Los cambios se aplican inmediatamente
\end{enumerate}

\section{Notificaciones}

El sistema de notificaciones le mantendrá informado sobre:

\begin{itemize}
    \item Actualizaciones disponibles
    \item Alertas de seguridad
    \item Tareas programadas
    \item Mensajes del sistema
\end{itemize}

\subsection{Centro de Notificaciones}

Acceda al centro de notificaciones haciendo clic en el icono \faIcon{bell} en la barra superior. Desde allí podrá:

\begin{itemize}
    \item Ver todas las notificaciones
    \item Marcar como leídas
    \item Configurar preferencias de notificación
    \item Acceder a acciones relacionadas
\end{itemize}

\notacientifica{Configure las notificaciones según sus preferencias para evitar sobrecarga de información. Puede seleccionar qué tipos de eventos desea recibir y la frecuencia de las alertas.}
